\begin{textsample}{POS Dim 1 – 2016 – Score 21.00 – 2016/t001203.txt}  \label{ex:f1_pos_052}
Here ' s a question that matters. [ Is it ethical to \textbf{evolve} \textbf{the} human body? ] Because we ' re beginning to get all \textbf{the} tools together to \textbf{evolve} ourselves.
And we can \textbf{evolve} bacteria and we can \textbf{evolve} plants and we can \textbf{evolve} animals, and we ' re now reaching a point where we really have to ask, is it really ethical and do we want to \textbf{evolve} human beings?
And as you ' re thinking about that, let me talk about that in \textbf{the} context of prosthetics, prosthetics past, present, future.
So this is \textbf{the} \textbf{iron} hand that belonged to one of \textbf{the} German counts.
Loved to fight, lost his arm in one of these battles.
No problem, he just made a suit of armor, put it on, perfect prosthetic.
That ' s where \textbf{the} concept of ruling with an \textbf{iron} fist comes from.
And of course these prosthetics have been getting more and more useful, more and more modern.
You can hold soft-boiled eggs.
You can have all types of controls, and as you ' re thinking about that, there are wonderful people like Hugh Herr who have been building absolutely extraordinary prosthetics.
So \textbf{the} wonderful Aimee Mullins will go out and say, how tall do I want to be tonight?
Or Hugh will say what type of cliff do I want to climb?
Or does somebody want to run a marathon, or does somebody want to ballroom dance?
And as you adapt these things, \textbf{the} interesting thing about prosthetics is they ' ve been coming inside \textbf{the} body.
So these external prosthetics have now become \textbf{artificial} knees.
They ' ve become \textbf{artificial} hips.
And then they ' ve \textbf{evolved} further to become not just nice to have but essential to have.
So when you ' re talking about a heart pacemaker as a prosthetic, you ' re talking about something that isn ' t just,"I ' m missing my leg,"it ' s,"if I don ' t have this, I can die."And at that point, a prosthetic becomes a symbiotic relationship with \textbf{the} human body.
And four of \textbf{the} smartest people that I ' ve ever met — Ed Boyden, Hugh Herr, Joe Jacobson, Bob Lander — are working on a Center for Extreme Bionics.
And \textbf{the} interesting thing of what you ' re seeing here is these prosthetics now get integrated into \textbf{the} bone.
They get integrated into \textbf{the} skin.
They get integrated into \textbf{the} muscle.
And one of \textbf{the} other sides of Ed is he ' s been thinking about how to connect \textbf{the} brain using light or other mechanisms directly to things like these prosthetics.
And if you can do that, then you can begin changing fundamental aspects of humanity.
So how quickly you react to something depends on \textbf{the} \textbf{diameter} of a nerve.
And of course, if you have nerves that are external or prosthetic, say with light or liquid metal, then you can increase that \textbf{diameter} and you could even increase it theoretically to \textbf{the} point where, as long as you could see \textbf{the} muzzle flash, you could step out of \textbf{the} way of a \textbf{bullet}.
Those are \textbf{the} order of magnitude of changes you ' re talking about.
This is a fourth sort of level of prosthetics.
These are Phonak hearing aids, and \textbf{the} reason why these are so interesting is because they cross \textbf{the} threshold from where prosthetics are something for somebody who is"disabled"and they become something that somebody who is"normal"might want to actually have, because what this prosthetic does, which is really interesting, is not only does it help you hear, you can focus your hearing, so it can hear \textbf{the} conversation going on over there.
You can have superhearing.
You can have hearing in 360 degrees.
You can have white noise.
You can record, and oh, by \textbf{the} way, they also put a phone into this.
So this functions as your hearing aid and also as your phone.
And at that point, somebody might actually want to have a prosthetic voluntarily.
All of these thousands of loosely connected little pieces are coming together, and it ' s about time we ask \textbf{the} question, how do we want to \textbf{evolve} human beings over \textbf{the} next century or two?
And for that we turn to a great philosopher who was a very smart man despite being a Yankee fan.
And Yogi Berra used to say, of course, that it ' s very tough to make predictions, especially about \textbf{the} future.
So instead of making a prediction about \textbf{the} future to begin with, let ' s take what ' s happening in \textbf{the} present with people like Tony Atala, who is redesigning 30-some-odd organs.
And maybe \textbf{the} ultimate prosthetic isn ' t having something external, titanium.
Maybe \textbf{the} ultimate prosthetic is take your own gene code, remake your own body parts, because that ' s a whole lot more effective than any kind of a prosthetic.
But while you ' re at it, then you can take \textbf{the} work of Craig Venter and Ham Smith.
And one of \textbf{the} things that we ' ve been doing is trying to figure out how to reprogram cells.
And if you can reprogram a cell, then you can change \textbf{the} cells in those organs.
So if you can change \textbf{the} cells in those organs, maybe you make those organs more radiation-resistant.
Maybe you make them absorb more oxygen.
Maybe you make them more efficient to filter out stuff that you don ' t want in your body.
And over \textbf{the} last few weeks, George Church has been in \textbf{the} news a lot because he ' s been talking about taking one of these programmable cells and inserting an entire human \textbf{genome} into that cell.
And once you can insert an entire human \textbf{genome} into a cell, then you begin to ask \textbf{the} question, would you want to enhance any of that \textbf{genome}?
Do you want to enhance a human body?
How would you want to enhance a human body?
Where is it ethical to enhance a human body and where is it not ethical to enhance a human body?
And all of a sudden, what we ' re doing is we ' ve got this multidimensional chess board where we can change human genetics by using viruses to attack things like AIDS, or we can change \textbf{the} gene code through gene therapy to do away with some hereditary diseases, or we can change \textbf{the} environment, and change \textbf{the} expression of those genes in \textbf{the} epigenome and pass that on to \textbf{the} next generations.
And all of a sudden, it ' s not just one little bit, it ' s all these \textbf{stacked} little bits that allow you to take little portions of it until all \textbf{the} portions coming together lead you to something that ' s very different.
And a lot of people are very scared by this stuff.
And it does sound scary, and there are risks to this stuff.
So why in \textbf{the} world would you ever want to do this stuff?
Why would we really want to alter \textbf{the} human body in a fundamental way? \textbf{The} answer lies in part with Lord Rees, astronomer royal of Great Britain.
And one of his favorite sayings is \textbf{the} \textbf{universe} is 100 percent malevolent.
So what does that mean?
It means if you take any one of your bodies at random, drop it anywhere in \textbf{the} \textbf{universe}, drop it in space, you die.
Drop it on \textbf{the} Sun, you die.
Drop it on \textbf{the} \textbf{surface} of Mercury, you die.
Drop it near a supernova, you die.
But fortunately, it ' s only about 80 percent effective.
So as a great physicist once said, there ' s these little upstream eddies of \textbf{biology} that create order in this rapid torrent of entropy.
So as \textbf{the} \textbf{universe} dissipates energy, there ' s these upstream eddies that create biological order.
Now, \textbf{the} problem with eddies is, they tend to \textbf{disappear}.
They shift.
They move in rivers.
And because of that, when an eddy shifts, when \textbf{the} Earth becomes a snowball, when \textbf{the} Earth becomes very hot, when \textbf{the} Earth gets hit by an \textbf{asteroid}, when you have supervolcanoes, when you have solar flares, when you have potentially extinction-level events like \textbf{the} next election — then all of a sudden, you can have periodic extinctions.
And by \textbf{the} way, that ' s happened five times on Earth, and therefore it is very likely that \textbf{the} human species on Earth is going to go extinct someday.
Not next week, not next month, maybe in November, but maybe 10, 000 years after that.
As you ' re thinking of \textbf{the} consequence of that, if you believe that extinctions are common and natural and normal and \textbf{occur} periodically, it becomes a moral imperative to diversify our species.
And it becomes a moral imperative because it ' s going to be really hard to live on Mars if we don ' t fundamentally modify \textbf{the} human body.
Right?
You go from one cell, mom and dad coming together to make one cell, in a cascade to 10 trillion cells.
We don ' t know, if you change \textbf{the} gravity substantially, if \textbf{the} same thing will happen to create your body.
We do know that if you expose our bodies as they currently are to a lot of radiation, we will die.
So as you ' re thinking of that, you have to really redesign things just to get to Mars.
Forget about \textbf{the} moons of Neptune or Jupiter.
And to borrow from Nikolai Kardashev, let ' s think about life in a series of scales.
So Life One civilization is a civilization that begins to alter his or her looks.
And we ' ve been doing that for thousands of years.
You ' ve got tummy tucks and you ' ve got this and you ' ve got that.
You alter your looks, and I ' m told that not all of those alterations take place for medical reasons.
Seems odd.
A Life Two civilization is a different civilization.
A Life Two civilization alters fundamental aspects of \textbf{the} body.
So you put human growth \textbf{hormone} in, \textbf{the} person grows taller, or you put x in and \textbf{the} person gets \textbf{fatter} or loses metabolism or does a whole series of things, but you ' re altering \textbf{the} functions in a fundamental way.
To become an intrasolar civilization, we ' re going to have to create a Life Three civilization, and that looks very different from what we ' ve got here.
Maybe you splice in Deinococcus radiodurans so that \textbf{the} cells can resplice after a lot of exposure to radiation.
Maybe you breathe by having oxygen flow through your \textbf{blood} instead of through your lungs.
But you ' re talking about really radical redesigns, and one of \textbf{the} interesting things that ' s happened in \textbf{the} last decade is we ' ve discovered a whole lot of \textbf{planets} out there.
And some of them may be Earth-like. \textbf{The} problem is, if we ever want to get to these \textbf{planets}, \textbf{the} fastest human \textbf{objects} — Juno and Voyager and \textbf{the} rest of this stuff — take tens of thousands of years to get from here to \textbf{the} nearest solar system.
So if you want to start exploring beaches somewhere else, or you want to see two-sun sunsets, then you ' re talking about something that is very different, because you have to change \textbf{the} timescale and \textbf{the} body of humans in ways which may be absolutely unrecognizable.
And that ' s a Life Four civilization.
Now, we can ' t even begin to imagine what that might look like, but we ' re beginning to get glimpses of instruments that might take us even that far.
And let me give you two examples.
So this is \textbf{the} wonderful Floyd Romesberg, and one of \textbf{the} things that Floyd ' s been doing is he ' s been playing with \textbf{the} basic \textbf{chemistry} of life.
So all life on this \textbf{planet} is made in ATCGs, \textbf{the} four letters of DNA.
All bacteria, all plants, all animals, all humans, all cows, everything else.
And what Floyd did is he changed out two of those base pairs, so it ' s ATXY.
And that means that you now have a parallel system to make life, to make babies, to reproduce, to \textbf{evolve}, that doesn ' t \textbf{mate} with most things on Earth or in fact maybe with nothing on Earth.
Maybe you make plants that are immune to all bacteria.
Maybe you make plants that are immune to all viruses.
But why is that so interesting?
It means that we are not a unique solution.
It means you can create alternate \textbf{chemistries} to us that could be \textbf{chemistries} adaptable to a very different \textbf{planet} that could create life and heredity. \textbf{The} second experiment, or \textbf{the} other implication of this experiment, is that all of you, all life is based on 20 amino acids.
If you don ' t substitute two amino acids, if you don ' t say ATXY, if you say ATCG + XY, then you go from 20 building blocks to 172, and all of a sudden you ' ve got 172 building blocks of amino acids to build life-forms in very different shapes. \textbf{The} second experiment to think about is a really weird experiment that ' s been taking place in China.
So this guy has been transplanting hundreds of mouse heads.
Right?
And why is that an interesting experiment?
Well, think of \textbf{the} first heart transplants.
One of \textbf{the} things they used to do is they used to bring in \textbf{the} wife or \textbf{the} daughter of \textbf{the} donor so \textbf{the} donee could tell \textbf{the} doctors,"Do you recognize this person?
Do you love this person?
Do you feel anything for this person?"We laugh about that today.
We laugh because we know \textbf{the} heart is a muscle, but for hundreds of thousands of years, or tens of thousands of years,"I gave her my heart.
She took my heart.
She broke my heart."We thought this was emotion and we thought maybe emotions were transplanted with \textbf{the} heart.
Nope.
So how about \textbf{the} brain?
Two possible outcomes to this experiment.
If you can get a mouse that is functional, then you can see, is \textbf{the} new brain a blank slate?
And boy, does that have implications.
Second option : \textbf{the} new mouse recognizes Minnie Mouse. \textbf{The} new mouse remembers what it ' s afraid of, remembers how to navigate \textbf{the} maze, and if that is true, then you can transplant memory and consciousness.
And then \textbf{the} really interesting question is, if you can transplant this, is \textbf{the} only input-output mechanism this down here?
Or could you transplant that consciousness into something that would be very different, that would last in space, that would last tens of thousands of years, that would be a completely redesigned body that could hold consciousness for a long, long period of time?
And let ' s come back to \textbf{the} first question : Why would you ever want to do that?
Well, I ' ll tell you why.
Because this is \textbf{the} ultimate selfie.
This is taken from six billion miles away, and that ' s Earth.
And that ' s all of us.
And if that little thing goes, all of humanity goes.
And \textbf{the} reason you want to alter \textbf{the} human body is because you eventually want a picture that says, that ' s us, and that ' s us, and that ' s us, because that ' s \textbf{the} way humanity survives long-term extinction.
And that ' s \textbf{the} reason why it turns out it ' s actually unethical not to \textbf{evolve} \textbf{the} human body even though it can be scary, even though it can be challenging, but it ' s what ' s going to allow us to explore, live and get to places we can ' t even dream of today, but which our great-great-great-great- grandchildren might someday.
Thank you very much.

% matched lemmas: artificial, asteroid, biology, blood, bullet, chemistry, diameter, disappear, evolve, fat, genome, hormone, iron, mate, object, occur, planet, stack, surface, the, universe
\end{textsample}
